\documentclass[11pt]{article}
\usepackage[a4paper,margin=2.4cm]{geometry}
\usepackage{fontspec}
\usepackage{xeCJK}
\setmainfont{Noto Serif}
\setCJKmainfont{Yu Gothic}
\setCJKsansfont{Yu Gothic}
\setlength{\parindent}{1em}
\setlength{\parskip}{0.5em}
\setlength{\emergencystretch}{2em}
\begin{document}

\section*{28. 群指標とイデアル論}

\begin{center}
J. Ber. d. DMV 34 (1925), S. 144
\end{center}

\textbf{3. E. Noether, ゲッティンゲン：群指標とイデアル論。}

Frobenius の群指標の理論、すなわち有限群の表現論は、完全可約環（今日の半単純環）である群環のイデアル論として捉えられる。したがって、完全可約環の加法的分解における同型定理が一般に展開される。直和としての、直既約な単純イデアルへの表示は同型を除いて一意であり、同じ両側イデアルに属する直既約片側イデアルは互いに同型である。したがって同型なイデアルを一つの類にまとめると、両側直既約イデアルの数と同じだけの異なる直既約イデアル類が存在する。

完全可約な超複素量の系、特に群環へ特殊化すると、直既約イデアル類と既約表現類は互いに一意に対応する。さらに両側直既約イデアルと指標が対応し、両側イデアルの加法的直既約成分の数は直既約イデアルの階数に等しい。これにより Frobenius の理論は位置づけられる。

詳しい叙述は \emph{Math. Ann.} に現れる予定である。

\end{document}
